\chapter{Project Overview}

\section{Background}
Human activity recognition (HAR) at home supports elderly-care monitoring,
rehabilitation and smart-home automation. Camera-based HAR is accurate, but it
records identifiable images of people in private spaces, it fails in darkness,
and many users do not accept it. Wearable sensors avoid images but must be worn
and charged. Single-chip frequency-modulated continuous-wave (FMCW) mmWave radars
operating at 60--81\,GHz (e.g.\ the Texas Instruments IWR1443/IWR1843 used by
both datasets in this project) sense through darkness and light clothing. For
each frame they output a small \emph{point cloud}: a set of detections with a 3-D
position, a radial (Doppler) velocity and an intensity. The output contains no
image and no visual identity, so the sensing is privacy-preserving by
construction.

\section{Problem statement}
Given a short sequence of radar point-cloud frames, the task is to classify the
human activity being performed, accurately and at a computational cost that
suits an edge device. The difficulty has four sources:
\begin{itemize}[nosep]
\item \textbf{Sparsity}: a frame holds only about 6 (MiliPoint) to 21 (RadHAR) returns
  (measured in Chapter~\ref{ch:data}).
\item \textbf{Noise}: many returns are antenna self-reflections, static clutter
  or multipath rather than the person. Our analysis shows that about half of all
  RadHAR returns are not on the person (Chapter~\ref{ch:pre}).
\item \textbf{Temporal nature}: the activity shows up in how the sparse cloud moves
  over time, not in the shape of a single frame.
\item \textbf{Evaluation leakage}: samples are overlapping windows cut from long
  recordings, so careless splitting puts near-identical samples in both the
  training and the test data (Chapter~\ref{ch:leak}).
\end{itemize}

\section{Motivation}
Published results on the same public radar datasets vary enormously. On
MiliPoint's 49-action task, the dataset authors report at most 34.45\% top-1
\cite{cui2023milipoint}, while a 2025 graph-network paper reports 98.25\% under a
random 80/10/10 split \cite{zakka2025makgcn}. A model that is accurate only
because of leakage is useless in a real home. This project therefore aims to
show what a lightweight model achieves under \emph{honest} evaluation and how
robust it is to realistic radar disturbances.

\section{Objectives}
\begin{enumerate}[nosep]
\item Build a reproducible pipeline for two public mmWave point-cloud datasets
  (MiliPoint as the primary dataset, RadHAR as the secondary one).
\item Design radar-specific pre-processing that removes non-human returns without
  removing human motion.
\item Develop a lightweight point-transformer classifier suitable for edge
  deployment and compare it with established baselines (PointNet, PointNet++,
  Point Transformer, DGCNN, PointNet+BiLSTM and the RadHAR voxel CNN-LSTM).
\item Evaluate under leakage-free protocols with three seeds, paired significance
  tests, ablations and a dedicated noise-robustness study.
\item Deliver a saved model, an inference script and a working interactive
  demonstration, together with a research paper, this report and a presentation.
\end{enumerate}

\section{Scope}
In scope: activity classification of single-person radar point-cloud sequences
using the label sets of RadHAR (5 activities) and MiliPoint (49 actions);
CPU-based efficiency measurement; robustness to synthetic perturbations of
recorded data.
Out of scope: fall detection (neither dataset has fall labels), multi-person
scenes, on-device hardware deployment, and clinical validation. No claim of
clinical suitability is made.

\section{Expected outcomes}
A trained and saved model with measured accuracy on held-out data; quantified
robustness under controlled noise; an automated leakage audit; a demonstration
application that runs the real model; and complete documentation.

\section{System requirements}
\begin{itemize}[nosep]
\item Python $\geq$ 3.10, PyTorch $\geq$ 2.1 (developed with 2.15), NumPy, SciPy,
  scikit-learn, pandas, Matplotlib, PyYAML; Streamlit and Plotly for the demo.
\item No GPU required. All experiments in this report ran on a 2-core CPU with
  7\,GB RAM (Section~\ref{sec:env}).
\item Disk: about 1\,GB for the RadHAR raw data and 0.4\,GB for MiliPoint.
\end{itemize}
